
 We represent the polyhedron
\[
P = \left\{
\begin{pmatrix}
x_1\\
x_2
\end{pmatrix}\in \mathbb{R}^2
| x_1 + x_2 \leq 1, 3 x_1 + 10 x_2 \leq 15, x_1 \geq 0, x_2 \geq 0 \right\}
\]
in the following figure:
\begin{center}
\begin{tikzpicture}[scale=1.9]
  \begin{axis}[ xlabel={$x_1$}, ylabel={$x_2$}
      ,axis on top=true
      ,axis equal
      ,axis lines=middle
      ,samples=41
      ,xtick = {-1,1}
      ,ytick = {1}
      ,domain=0:1.5
      ,xmin=-1,xmax=5.2
      ,ymin=-0.3,ymax=2
,legend pos=south east
]
\addplot [thick,color=gray!20,fill=gray!20, 
                    fill opacity=0.05]coordinates {
            (0, 1) 
            (1, 0)
            (0, 0)  };
\addplot[color=black, thin,domain=-1:1.2] {1-x}
node[pos=0.5,sloped,above] {\tiny $x_1+x_2=1$};
\addplot[color=black, thin,domain=-1:5.1] {-0.3 * x + 1.5}
node[pos=0.5,sloped,above] {\tiny $3x_1+10x_2=15$};
\node [circle,fill=black,draw,scale=0.2,label=below left:$A$] at (0,0) {A};
\node [circle,fill=black,draw,scale=0.2,label=right:$B$] at (1,0) {B};
\node [circle,fill=black,draw,scale=0.2,label=below:$C$] at (0,1) {C};
\node [circle,fill=black,draw,scale=0.2,label=below:$D$] at (-5/7,12/7) {D};
\node [circle,fill=black,draw,scale=0.2,label=below:$E$] at (5,0) {E};
\node [circle,fill=black,draw,scale=0.2,label=above right:$F$] at (0,1.5) {F};

\end{axis}
\end{tikzpicture}
\end{center}


The corresponding polyhedron in standard form is
$$Q = \left\{
\begin{pmatrix}
x_1\\
x_2\\
x_3\\
x_4
\end{pmatrix} \in \R^4
| Ax = b, x \geq 0 \right\},$$
with
$$A =
\begin{pmatrix*}[r]
1 & 1 & 1 & 0\\
3 & 10 & 0 & 1
\end{pmatrix*},
b = 
\begin{pmatrix}
1 \\
15
\end{pmatrix},$$
with the variables $x_3$ and $x_4$ being slack variables.
Each basic solution is obtained by selecting 2 variables out of 4 to
be in the basis. There is a total of 6 possible selections of basic variables.\\
The matrix $B=(A_{j1},...,A_{jm})$ is composed of columns $j_1,...,j_m$ of the matrix $A$.


\begin{enumerate}

\item Basic solution with $x_1$ and $x_2$ in the basis ($j_1=1$,$j_2=2$).
\[
B = 
\begin{pmatrix*}[r]
1 & 1\\
3 &10
\end{pmatrix*};
B^{-1} =
\begin{pmatrix*}[r]
\frac{10}{7} & -\frac{1}{7}\\
-\frac{3}{7} &\frac{1}{7}
\end{pmatrix*};
x_B = B^{-1}b =
\begin{pmatrix}
-\frac{5}{7} \\
\frac{12}{7}
\end{pmatrix};
x =
\begin{pmatrix}
-\frac{5}{7} \\
\frac{12}{7} \\
0 \\
0
\end{pmatrix}.
\]

This basic solution is not feasible because $B^{-1}b\ngeqslant0$. It corresponds to point $D=(-\frac{5}{7},\frac{12}{7})$ in the figure.

\item Basic solution with $x_1$ and $x_3$ in the basis $(j_1=1,j_2=3)$.
\[
B = 
\begin{pmatrix*}[r]
1 & 1\\
3 &0
\end{pmatrix*};
B^{-1} =
\begin{pmatrix*}
0 & \frac{1}{3}\\
1 & -\frac{1}{3}
\end{pmatrix*};
x_B = B^{-1}b =
\begin{pmatrix}
5 \\
-4
\end{pmatrix};
x =
\begin{pmatrix}
5 \\
0 \\
-4 \\
0
\end{pmatrix}.
\]

This basic solution is not feasible because $B^{-1}b\ngeqslant0$. It corresponds to point $E=(5,0)$.

\item Basic solution with $x_1$ and $x_4$ in the basis $(j_1=1,j_2=4)$.
\[
B = 
\begin{pmatrix}
1 & 0\\
3 &1
\end{pmatrix};
B^{-1} =
\begin{pmatrix*}
1 & 0\\
-3 & 1
\end{pmatrix*};
x_B = B^{-1}b =
\begin{pmatrix}
1 \\
12
\end{pmatrix};
x =
\begin{pmatrix}
1 \\
0 \\
0 \\
12
\end{pmatrix}.
\]

This basic solution is feasible and corresponds to point $B=(1,0)$ .

\item Basic solution with $x_2$ and $x_3$ in the basis $(j_1=2,j_2=3)$.
\[
B =
\begin{pmatrix*}[r]
1 & 1\\
10 &0
\end{pmatrix*};
B^{-1} =
\begin{pmatrix*}[r]
0 & \frac{1}{10}\\
1 & -\frac{1}{10}
\end{pmatrix*};
x_B = B^{-1}b =
\begin{pmatrix*}[r]
\frac{3}{2} \\
-\frac{1}{2}
\end{pmatrix*};
x =
\begin{pmatrix*}[r]
0 \\
\frac{3}{2} \\
-\frac{1}{2} \\
0
\end{pmatrix*}.
\]

This basic solution is not feasible because $B^{-1}b\ngeqslant0$. It corresponds to point $F=(0,\frac{3}{2})$ in the figure.

\item Basic solution with $x_2$ and $x_4$ in the basis $(j_1=2,j_2=4)$.
\[
B = 
\begin{pmatrix*}[r]
1 & 0\\
10 &1
\end{pmatrix*};
B^{-1} =
\begin{pmatrix}
1 & 0\\
-10 & 1
\end{pmatrix};
x_B = B^{-1}b =
\begin{pmatrix}
1 \\
5
\end{pmatrix};
x =
\begin{pmatrix}
0 \\
1 \\
0 \\
5
\end{pmatrix}.
\]

This basic solution is feasible and corresponds to point $C=(0,1)$ in the figure.

\item Basic solution with $x_3$ and $x_4$ in the basis $(j_1=3,j_2=4)$.
\[
B = B^{-1} =
\begin{pmatrix}
1 & 0\\
0 & 1
\end{pmatrix};
x_B = B^{-1}b =
\begin{pmatrix}
1 \\
15
\end{pmatrix};
x =
\begin{pmatrix}
0 \\
0 \\
1 \\
15
\end{pmatrix}.
\]

This basic solution is feasible and corresponds to point $A=(0,0)$ in the figure.
\end{enumerate}

In summary, out of the six basic solutions, only three are feasible
and correspond to vertices of the constraint polyhedron. 


